%------------------------------------------------------------------------------%
\documentclass[fleqn,utf8,aspectratio=169,12pt]{beamer}

\usepackage{gaal}

\title{Posições Relativas de Retas}

%------------------------------------------------------------------------------%
\begin{document}

\addtitle

%------------------------------------------------------------------------------%
\section{Posições Relativas no Plano}

%------------------------------------------------------------------------------%
\begin{frame}{Posições Relativas no Plano}
  \centering%
  \includegraphics{F-02-Posicoes_relativas_de_retas-conc-1}\;\pause%
  \includegraphics{F-02-Posicoes_relativas_de_retas-para-1}\;\pause%
  \includegraphics{F-02-Posicoes_relativas_de_retas-coin-1}%
\end{frame}

%------------------------------------------------------------------------------%
\begin{frame}{Posições Relativas entre Duas Retas no Plano}
  No plano, duas retas podem ser
  \pause
  \begin{itemize}
      \item concorrentes        \pause \tabto{36mm} --- \quad se tocam em um único ponto \pause
      \item paralelas distintas \pause \tabto{36mm} --- \quad não se tocam               \pause
      \item coincidentes        \pause \tabto{36mm} --- \quad são a mesma reta
  \end{itemize}
\end{frame}

%------------------------------------------------------------------------------%
\begin{frame}{Determinando a Posição Relativa entre Duas Retas no Plano}
  Dadas duas retas
  \begin{align*}
    r\colon \; \vec{r} & = \vec{p} + t    \, \vec{v} \\
    s\colon \; \vec{r} & = \vec{q} + \tau \, \vec{w}
  \end{align*}

  \pause
  Analisamos os vetores diretores  \(\vec{v}\) e \(\vec{w}\)
\end{frame}

%------------------------------------------------------------------------------%
\begin{frame}<handout:0>{Posições Relativas}
  \centering%
  \includegraphics{F-02-Posicoes_relativas_de_retas-conc-1}\;%
  \includegraphics{F-02-Posicoes_relativas_de_retas-para-1}\;%
  \includegraphics{F-02-Posicoes_relativas_de_retas-coin-1}%
\end{frame}

%------------------------------------------------------------------------------%
\begin{frame}<handout:0>{Posições Relativas}
  \centering%
  \includegraphics{F-02-Posicoes_relativas_de_retas-conc-2}\;%
  \includegraphics{F-02-Posicoes_relativas_de_retas-para-1}\;%
  \includegraphics{F-02-Posicoes_relativas_de_retas-coin-1}%
\end{frame}

%------------------------------------------------------------------------------%
\begin{frame}<handout:0>{Posições Relativas}
  \centering%
  \includegraphics{F-02-Posicoes_relativas_de_retas-conc-2}\;%
  \includegraphics{F-02-Posicoes_relativas_de_retas-para-2}\;%
  \includegraphics{F-02-Posicoes_relativas_de_retas-coin-1}%
\end{frame}

%------------------------------------------------------------------------------%
\begin{frame}{Posições Relativas}
  \centering%
  \includegraphics{F-02-Posicoes_relativas_de_retas-conc-2}\;%
  \includegraphics{F-02-Posicoes_relativas_de_retas-para-2}\;%
  \includegraphics{F-02-Posicoes_relativas_de_retas-coin-2}%
\end{frame}

%------------------------------------------------------------------------------%
\begin{frame}{Retas Paralelas no Plano}
  Se $\vec{v}$ e $\vec{w}$ \emph{não} são paralelos,
  \pause
  as retas são \emph{concorrentes}

  \bigskip
  \pause
  Se os vetores são \emph{paralelos},
  \pause
  i.e., existe $\alpha\in\R$ tal que
  \[
    \vec{w} = \alpha\vec{v}
  \]
  \pause
  as retas são
  \pause
  \begin{itemize}
    \item \emph{coincidentes}, \pause se possuem um ponto em comum \pause
    \item \emph{paralelas distintas}, caso contrário
  \end{itemize}
\end{frame}

%------------------------------------------------------------------------------%
%------------------------------------------------------------------------------%
\begin{question}
  Determine a posição relativa entre as retas $r$ e $s$
  \[
    r\colon
    \begin{cases}
      x = 1+2t \\
      y = 3-t
    \end{cases}
  \]
  \[
    s\colon
    \begin{cases}
      x = -2 + 4\tau \\
      y =  5 - 2\tau
    \end{cases}
  \]

\end{question}

%------------------------------------------------------------------------------%
\begin{resolution}{Solução}
  Os vetores diretores são
  \[
    \pause
    \vec v=(2,-1)
  \]
  \[
    \pause
    \vec w=(4,-2)
  \]

  \pause
  Como
  \[
    \vec w = 2\vec v
  \]
  \pause
  os vetores são paralelos

  \pause
  Então as retas são paralelas distintas ou coincidentes
\end{resolution}

%------------------------------------------------------------------------------%
\begin{resolution}{Solução}
  Igualando as coordenadas,
  \pause
  temos o sistema
  \[
    \begin{cases}
      1+2t = -2+4\tau \\
      3-t  = 5-2\tau
    \end{cases}
  \]
  \pause
  \[
    \begin{cases}
      2t -4\tau = -3 \\
      -t +2\tau =  2
    \end{cases}
  \]
  \pause
  \[
    \begin{cases}
      2t -4\tau = -3 \\
      2t -4\tau = -4
    \end{cases}
  \]
  \pause
  As equações são incompatíveis,
  \pause
  portanto, as retas são paralelas distintas
\end{resolution}

%------------------------------------------------------------------------------%
\begin{resolution}{Solução Alternativa}
  Se as retas são coincidentes todos os pontos em $r$ também estão em $s$

  \pause
  Escolhemos o ponto associado a $t=0$ na reta $r$
  \[
    \pause
    (x, y) = (1, 3)
  \]
  \pause
  Subistituindo nas equações da reta $s$ temos
\vspace{-\baselineskip}
\begin{columns}[t]
\column{0.4\textwidth}
  \[
    \pause
    \begin{cases}
      1 = -2 + 4\tau \\
      3 =  5 - 2\tau
    \end{cases}
  \]
  \[
    \pause
    \begin{cases}
      4\tau = 3 \\
      2\tau = -2
    \end{cases}
  \]
\column{0.6\textwidth}
  \[
    \pause
    \begin{cases}
      \tau = \sfrac{3}{4} \\
      \tau = \sfrac{-2}{2}
    \end{cases}
  \]
  \pause
  As equações são incompatíveis, \\
  \pause
  portanto, as retas são paralelas distintas
\end{columns}
\end{resolution}

%------------------------------------------------------------------------------%
\endinput

%------------------------------------------------------------------------------%
\begin{question}
  Determine a posição relativa e o ponto de interseção, caso exista,
  \[
    r\colon
    \begin{cases}
      x = 1 + t  \\
      y = 2 +2t
    \end{cases}
  \]
  \[
    s\colon
    \begin{cases}
      x = 4 - \tau  \\
      y = 8 -3\tau
    \end{cases}
  \]
\end{question}

%------------------------------------------------------------------------------%
\begin{resolution}{Solução}
  Os vetores diretores são
  \[
    \vec v=(1,2)
  \]
  \[
    \vec w=(-1,-3)
  \]

  \pause
  Eles não são paralelos

  \pause
  Portanto, no plano, as retas são concorrentes
\end{resolution}

%------------------------------------------------------------------------------%
\begin{resolution}{Ponto de Interseção}
\begin{columns}
\column{0.5\textwidth}
  \[\onslide<+->{
    \begin{cases}
      1 +  t = 4 -  \tau \\
      2 + 2t = 8 - 3\tau
    \end{cases}}
  \]
  \[\onslide<+->{
    \begin{cases}
       t = 3 - \tau \\
      2t + 3\tau = 6
    \end{cases}}
  \]
  \begin{align*}
    \onslide<+->{2(3 - \tau) + 3\tau & = 6 \\}
    \onslide<+->{6 - 2\tau + 3\tau   & = 6 \\}
    \onslide<+->{\tau                & = 0}
  \end{align*}
\column{0.5\textwidth}
  \[
   \onslide<+->{t = 3 - \tau}
   \onslide<+->{= 3}
  \]
  \onslide<+->{Substituindo na reta $s$}
  \begin{align*}
    \onslide<+->{x & = 4 - \tau}\onslide<+->{ = 4} \\
    \onslide<+->{y & = 8 -3\tau}\onslide<+->{ = 8}
  \end{align*}
  \onslide<+->{Assim
  \(
    r\cap s = \{(4, 8)\}
  \)}
\end{columns}
\end{resolution}

%------------------------------------------------------------------------------%
\endinput


%------------------------------------------------------------------------------%
\section{Posições Relativas no Espaço}

%------------------------------------------------------------------------------%
\begin{frame}{Posições Relativas entre Retas no Espaço}
  No \emph{espaço}, duas retas podem ser
  \pause
  \begin{itemize}
    \item concorrentes        \pause \tabto{34mm} --- \quad se tocam em um único ponto \pause
    \item paralelas distintas \pause \tabto{34mm} --- \quad ``mantêm a distância''     \pause
    \item coincidentes        \pause \tabto{34mm} --- \quad são a mesma reta           \pause
    \item reversas            \pause \tabto{34mm} --- \quad
          não são paralelas \pause
          e não possuem ponto de interseção
  \end{itemize}
\end{frame}

%------------------------------------------------------------------------------%
\begin{frame}{Retas Reversas em $\R^3$}
  Dadas duas retas
  \begin{align*}
    r\colon \; \vec{r} & = \vec{p} + t    \, \vec{v} \\
    s\colon \; \vec{r} & = \vec{q} + \tau \, \vec{w}
  \end{align*}
  \pause
  Verificamos se existe $\alpha$, tal que
  \[
    \vec{v} = \alpha\vec{w}
  \]
  \pause
  Se os vetores forem paralelos as retas são
  \emph{paralelas distintas} ou \emph{coincidentes}
\end{frame}

%------------------------------------------------------------------------------%
\begin{frame}{Retas Reversas em $\R^3$}
  Se os vetores \emph{não} forem paralelos

  \pause
  \bigskip
  Se as retas estiverem no mesmo plano elas são \emph{concorrentes}

  \pause
  \bigskip
  Se estiverem em planos diferentes são \emph{reversas}
\end{frame}

%------------------------------------------------------------------------------%
\begin{frame}<handout:0>{Retas Reversas}
  \centering%
  \includegraphics{F-02-Posicoes_relativas_de_retas-conc-2}\;%
  \includegraphics{F-02-Posicoes_relativas_de_retas-para-2}\;%
  \includegraphics{F-02-Posicoes_relativas_de_retas-coin-2}%
\end{frame}

%------------------------------------------------------------------------------%
\begin{frame}<handout:0>{Retas Reversas}
  \centering%
  \includegraphics{F-02-Posicoes_relativas_de_retas-conc-3}\;%
  \includegraphics{F-02-Posicoes_relativas_de_retas-para-2}\;%
  \includegraphics{F-02-Posicoes_relativas_de_retas-coin-2}%
\end{frame}

%------------------------------------------------------------------------------%
\begin{frame}<handout:0>{Retas Reversas}
  \centering%
  \includegraphics{F-02-Posicoes_relativas_de_retas-conc-3}\;%
  \includegraphics{F-02-Posicoes_relativas_de_retas-para-3}\;%
  \includegraphics{F-02-Posicoes_relativas_de_retas-coin-2}%
\end{frame}

%------------------------------------------------------------------------------%
\begin{frame}{Retas Reversas}
  \centering%
  \includegraphics{F-02-Posicoes_relativas_de_retas-conc-3}\;%
  \includegraphics{F-02-Posicoes_relativas_de_retas-para-3}\;%
  \includegraphics{F-02-Posicoes_relativas_de_retas-coin-3}%
\end{frame}

%------------------------------------------------------------------------------%
\begin{frame}{Retas Reversas em $\R^3$}
  As retas são coplanares,
  \pause
  se os vetores
  \(
    \vec{v}
  \),
  \(
    \vec{w}
  \)
  e
  \(
    \vec{q}-\vec{p}
  \)
  são coplanares
  \pause
  \[
    \left(\vec{v} \times \vec{w}\right) \cdot \left(\vec{p}-\vec{q}\right) = 0
  \]
  \pause
  nesse caso, as retas são \emph{concorrentes}

  \pause
  \bigskip
  Se
  \(
    \left(\vec{v} \times \vec{w}\right) \cdot \left(\vec{p}-\vec{q}\right) \neq 0
  \)
  \pause
  as retas são \emph{reversas}
\end{frame}

%------------------------------------------------------------------------------%
%------------------------------------------------------------------------------%
\begin{question}
  Determine a posição relativa entre as retas e, caso exista, o ponto de interseção
  \[
    r\colon
    \begin{cases}
      x = 1 +  t  \\
      y = 2 -  t  \\
      z = 3 + 2t
    \end{cases}
  \qquad\qquad
    s\colon
    \begin{cases}
      x = 3 - \tau \\
      y =     \tau \\
      z = 7 -2\tau
    \end{cases}
  \]
\end{question}

%------------------------------------------------------------------------------%
\begin{resolution}{Solução}
  Os vetores diretores são
  \[
    \vec{v} = (1, -1, 2)
  \]
  \[
    \vec{w} = (-1, 1, -2)
  \]

  \pause
  Como
  \[
    \vec{w} = -\vec{v}
  \]
  \pause
  as retas são paralelas
\end{resolution}

%------------------------------------------------------------------------------%
\begin{resolution}{Solução}
  Para verificar se são coincidentes,
  \pause
  comparamos os pontos
  \[
    P = (1, 2, 3)
  \]
  \[
    Q = (3, 0, 7)
  \]
  \pause
  Temos
  \[
    Q - P
    \pause
    = (2, -2, 4)
    \pause
    = 2\vec{v}
  \]
  \pause
  Portanto, o ponto $Q$ pertence à reta $r$

  \pause
  Consequentemente, as duas retas são coincidentes
\end{resolution}

%------------------------------------------------------------------------------%
\endinput

%------------------------------------------------------------------------------%
\begin{question}
  Determine a posição relativa entre as retas
  \[
    r\colon
    \begin{cases}
      x =  1 +  t \\
      y =  2 + 2t \\
      z = -1 + 3t
    \end{cases}
  \qquad\qquad
    s\colon
    \begin{cases}
      x & = 4 + 2\tau \\
      y & = 1 + 4\tau \\
      z & = 2 + 6\tau
    \end{cases}
  \]
\end{question}

%------------------------------------------------------------------------------%
\begin{resolution}{Solução}
\begin{columns}[t]
\column{0.3\textwidth}
  \onslide<+->{Os vetores diretores são}
  \[
    \onslide<+->{\vec{v} = (1, 2, 3)}
  \]
  \[
    \onslide<+->{\vec{w} = (2, 4, 6)}
  \]
  \onslide<+->{Como
  \[
    \vec{w} = 2\vec{v}
  \]}
  \onslide<+->{as retas são paralelas}
\column{0.3\textwidth}
  \onslide<+->{Verificar se
  \[
    P = (1, 2, -1) \in s
  \]}
  \vspace{-\baselineskip}
  \begin{align*}
    \onslide<+->{x     & = 4 + 2\tau    \\}
    \onslide<+->{1     & = 4 + 2\tau    \\}
    \onslide<+->{2\tau & = -3           \\}
    \onslide<+->{\tau  & = \frac{-3}{2}}
  \end{align*}
\column{0.3\textwidth}
  \begin{align*}
    \onslide<+->{y & = 1 + 4\tau         \\}
    \onslide<+->{  & = 1 + 4\frac{-3}{2} \\}
    \onslide<+->{  & = -5 \\}
    \onslide<+->{  & \neq 2}
  \end{align*}
  \onslide<+->{portanto, \(P\notin s\)}

  \medskip
  \onslide<+->{As retas são\\ paralelas distintas}
\end{columns}
\end{resolution}

%------------------------------------------------------------------------------%
\endinput

%------------------------------------------------------------------------------%
\begin{question}
  Determine a posição relativa entre as retas
  \[
    r\colon
    \begin{cases}
      x = t   \\
      y = 1+t \\
      z = 2t
    \end{cases}
  \qquad\qquad
    s\colon
    \begin{cases}
      x = 1+\tau \\
      y = 2-\tau \\
      z = 1+\tau
    \end{cases}
  \]
\end{question}

%------------------------------------------------------------------------------%
\begin{resolution}{Solução}
  Os vetores diretores são
  \[
    \vec{v} = (1, 1, 2)
  \]
  \[
    \vec{w} = (1, -1, 1)
  \]

  \pause
  Eles não são paralelos

  \pause
  Então, as retas podem ser concorrentes ou reversas
\end{resolution}

%------------------------------------------------------------------------------%
\begin{resolution}{Solução}
  Construindo o vetor conectando as duas retas
  \[
    \pause
    \vec{p} = (0, 1, 0)
  \]

  \[
    \pause
    \vec{q} = (1, 2, 1)
  \]

  \[
    \pause
    \vec{d}
    \pause
    = \vec{q} - \vec{p}
    \pause
    = (1, 2, 1) - (0, 1, 0)
    \pause
    = (1, 1, 1)
  \]

  \pause
  Precisamos verificar se os vetores $\vec{v}$, $\vec{w}$ e $\vec{d}$ são coplanares
\end{resolution}

%------------------------------------------------------------------------------%
\begin{resolution}{Solução}
  \onslide<+->{Produto misto}
  \begin{align*}
    \onslide<+->{(\vec{v} \times \vec{w}) \cdot \vec{d}}
    \onslide<+->{& = \begin{vmatrix}
          1 & 1  & 2 \\
          1 & -1 & 1 \\
          1 & 1  & 1
        \end{vmatrix}}
    \onslide<+->{\;=\; 1
      \begin{vmatrix}
        -1 & 1 \\
        1  & 1
      \end{vmatrix}
      - 1
      \begin{vmatrix}
        1 & 1 \\
        1 & 1
      \end{vmatrix}
      + 2
      \begin{vmatrix}
        1 & -1 \\
        1 & 1
      \end{vmatrix}}
    \\
    \onslide<+->{& = (-1-1)-0+2(1+1) \\}
    \onslide<+->{& = 2}
    \onslide<+->{\neq 0}
  \end{align*}

  \onslide<+->{Os vetores não são coplanares}

  \onslide<+->{Consequentemente, as retas são reversas}
\end{resolution}

%------------------------------------------------------------------------------%
\endinput


%------------------------------------------------------------------------------%
\section{Procedimento Geral}

%------------------------------------------------------------------------------%
\begin{frame}{Determinar a Posição Relativa}
  Dadas as retas
  \begin{align*}
    r\colon \; \vec{r} & = \vec{p} + t    \, \vec{v} \\
    s\colon \; \vec{r} & = \vec{q} + \tau \, \vec{w}
  \end{align*}

  \pause
  Verificamos se os vetores diretores são paralelos
  \[
    \vec{w} \;\stackrel{?}{=}\; \alpha \vec{v}
  \]
\end{frame}

%------------------------------------------------------------------------------%
\begin{frame}{Determinar a Posição Relativa}
  Se os vetores diretores forem \emph{paralelos}
  \[
    \pause
    \vec{v} = \alpha \vec{w}
  \]
  \pause
  as retas são paralelas ou coincidentes

  \bigskip
  \pause
  Se
  \(
    \vec{p}-\vec{q}
  \)
  \pause
  é paralelo a \(\vec{v}\)
  \pause
  e as retas são coincidentes
\end{frame}

%------------------------------------------------------------------------------%
\begin{frame}{Determinar a Posição Relativa}
  Se os vetores \emph{não} forem paralelos

  \bigskip
  \pause
  No \emph{plano}: as retas são concorrente

  \bigskip
  \pause
  No \emph{espaço}: calculamos o produto misto \;
  \(
    \left(\vec{v} \times \vec{w}\right) \cdot \left(\vec{p}-\vec{q}\right)
  \)

  \medskip
  \pause
  \qquad Se o produto misto for zero,
  \pause
  as retas são coplanares e concorrentes

  \pause
  \qquad Se o produto misto for diferente de zero,
  \pause
  as retas são reversas
\end{frame}

%------------------------------------------------------------------------------%
%%------------------------------------------------------------------------------%
\begin{question}
  Determine a posição relativa entre as retas
  \[
    r\colon
    \begin{cases}
      x = 2+t   \\
      y = -1+2t \\
      z = 1-t
    \end{cases}
  \qquad\qquad
    s\colon
    \begin{cases}
      x = 1+2\tau \\
      y = 1+4\tau \\
      z = 3-2\tau
    \end{cases}
  \]
\end{question}

%------------------------------------------------------------------------------%
\begin{resolution}{Solução}
  Os vetores diretores são
  \[
    \vec{v} = (1, 2, -1)
  \]
  \[
    \vec{w} = (2, 4, -2)
  \]
  \[
    \vec{w} = 2\vec{v}
  \]
  as retas são paralelas ou coincidentes
\end{resolution}

%------------------------------------------------------------------------------%
\begin{resolution}{Solução}
  Escolhemos
  \[
  P_0=(2,-1,1)
  \]
  e
  \[
  Q_0=(1,1,3)
  \]
  Então,
  \[
  Q_0-P_0=(-1,2,2)
  \]
  Se as retas fossem coincidentes, esse vetor deveria ser paralelo a
  \[
  \vec v=(1,2,-1)
  \]
  Isso não ocorre.

  De fato, comparando as componentes:
  \[
  \frac{-1}{1}\neq\frac{2}{2}
  \]
  Logo, as retas são \textbf{paralelas distintas}.
\end{resolution}

%------------------------------------------------------------------------------%
\endinput

%%------------------------------------------------------------------------------%
%\section{Lista Mínima}
%
%%------------------------------------------------------------------------------%
%\begin{frame}{Lista Mínima}
%  \begin{enumerate}
%    \item
%  \end{enumerate}
%
%  \vfill
%  Atenção: A prova é baseada no livro, não nas apresentações
%\end{frame}

%------------------------------------------------------------------------------%
\end{document}
%------------------------------------------------------------------------------%
